\subsection{\normalsize{Otros métodos de optimización de memoria}}
Además de la compresión sin pérdida, se han desarrollado enfoques alternativos
para reducir el consumo de memoria en simuladores cuánticos, incluyendo
técnicas basadas en redes tensoriales, poda de estados y optimización HPC.

Uno de los primeros enfoques en esta dirección fue propuesto por \textcite{Markov2008},
quienes introdujeron el uso de redes tensoriales para simular circuitos
cuánticos mediante contracción optimizada de tensores. En este método, el
circuito se representa como un grafo tensorial, y su simulación se optimiza contrayendo
secuencias de tensores en un orden eficiente. Esto permite representar ciertos
circuitos con memoria subexponencial, aunque la efectividad del método depende
de la estructura del entrelazamiento.

Posteriormente, \textcite{Boixo2018} propusieron un modelo gráfico para
circuitos cuánticos de baja profundidad basado en la eliminación de variables
en grafos probabilísticos. Su técnica permitió simular circuitos aleatorios cercanos
al régimen de supremacía cuántica en hardware clásico, extendiendo el tamaño
de circuitos simulables sin necesidad de almacenar el estado cuántico completo.
Sin embargo, este método pierde eficacia conforme aumenta la profundidad del
circuito.

En paralelo, \textcite{Pednault2017} introdujeron técnicas de diferimiento de contracciones
tensoriales en simulaciones HPC. Su propuesta permite manejar circuitos más
grandes al intercambiar memoria por cómputo adicional, almacenando resultados intermedios
en memoria secundaria y evitando la expansión total del estado cuántico en RAM.
Esta estrategia ha sido clave en la simulación de circuitos con estructuras
altamente no triviales, aunque su implementación requiere acceso a infraestructuras
de supercomputación.

Finalmente, como se analizó previamente, la tesis doctoral de \textcite{DiazToro2025}
introduce técnicas de compresión con pérdida de precisión en esquemas distribuidos
de simulación cuántica. Este antecedente resulta especialmente relevante porque
consolida el uso de bibliotecas científicas de compresión, vectores de estado y
paralelismo distribuido con \emph{MPI} como una línea válida para reducir consumo
de memoria en simulación cuántica. Al mismo tiempo, deja visible un espacio de
estudio complementario: evaluar hasta qué punto técnicas de compresión
\textbf{sin pérdida de precisión} pueden aportar reducciones útiles de memoria sin
introducir discrepancias numéricas en el estado simulado.
